\subsection{Surjective Functions}
\label{subsec:surjective-functions}

\begin{tcolorbox}[title={Definition --- Surjective}]
\begin{definition}[Surjective Function]
\label{def:alg-surjective}
Let $f:A\to B$.
The function $f$ is \textbf{surjective} if
\[
\forall b \in B,\; \exists a \in A,\; f(a)=b.
\]
\end{definition}
\end{tcolorbox}

\begin{remark*}[Standard quantified statement]
\[
\operatorname{Surjective}(f,A,B)
\Longleftrightarrow
f:A\to B
\land
\forall b\in B,\; \exists a\in A,\; f(a)=b.
\]
\end{remark*}

\begin{remark*}[Definition predicate reading]
$f$ is surjective exactly when every element of the codomain is hit by some
input.
\end{remark*}

\begin{remark*}[Negated quantified statement]
\[
\neg\operatorname{Surjective}(f,A,B)
\Longleftrightarrow
\exists b\in B,\; \forall a\in A,\; f(a)\neq b.
\]
\end{remark*}

\begin{remark*}[Interpretation]
Surjectivity says that the image of $f$ fills the whole codomain.
\end{remark*}

\begin{remark*}[Dependencies]
\begin{itemize}
  \item Function.
\end{itemize}
\end{remark*}
